\section{Methods}

% ---------------------------------------------------------------------------------------------
\subsection{Descriptor selection}
% ---------------------------------------------------------------------------------------------

The descriptor set is the union of RDKit's 217 \texttt{Descriptors.\_descList} entries and
Mordred's 1{,}613 two-dimensional descriptors, 1{,}830 columns in total. Mordred's 3D block is
excluded throughout: it requires conformer generation, which is two orders of magnitude more
expensive. Selection is based on a run of 20{,}000 ChEMBL molecules where columns that are
constant or less than 95\% finite are dropped first. A descriptor is then kept unless it
correlates $|\rho| \geq 0.99$ with one already kept, yielding \textbf{865 descriptors}.

% ---------------------------------------------------------------------------------------------
\subsection{Fingerprint}
% ---------------------------------------------------------------------------------------------

Morgan fingerprints are generated from RDKit's \texttt{rdFingerprintGenerator}, folded to 2048
bits, with \texttt{includeChirality=True} and radius $r=3$, which is slightly better at
resolving similar rings, the classical cyclopentane versus cyclohexane problem \cite{morgan}.
This does lead to slightly more hash collisions. Widening to 4096 bits halves the loss and
doubles the memory, which we judged not worth it.

% ---------------------------------------------------------------------------------------------
\subsection{Optimized descriptor implementations}
% ---------------------------------------------------------------------------------------------

Most descriptors are re-implemented in C++ in the obvious way, and only the departures are
described here. The specification is the reference source code, not its documentation, and three
upstream defects are reproduced rather than repaired, because a descriptor that silently
disagrees with the library everyone else uses is worse than one that is consistently odd: five
E-state SMARTS rows are written \texttt{[N,nD2H0]} without a semicolon, so the aliphatic
alternative carries no degree or hydrogen constraint at all, while the neighbouring carbon rows
use \texttt{[C,c;D2H0]} and do not; \texttt{[SeD2H0]} and its equivalents for Li, Be, Si, Ge, As,
Sn and Pb parse to element-\emph{number} queries with no aromaticity constraint, so
\texttt{aaSe} does fire on selenophene's aromatic selenium; and RDKit's shipped
\texttt{Data/Crippen.txt} is stale relative to the table compiled into the library, two of whose
110 rows negate by element rather than by aliphatic symbol, which we pin with behavioural probes
so that a future release fixing the text file leaves the port unaffected. The dividing test is
whether the upstream quantity is a function of the molecule at all; two are not, and we diverge
from them by design. Mordred's \texttt{InformationContent} mutates a visited set while iterating
over it, leaving roughly 20\% of molecules with an atom-equivalence code that depends on input
atom ordering, and \texttt{ExtendedTopochemicalAtom} scores a pyrrole nitrogen 0 where a pyridine
nitrogen scores 2, with the outcome turning on atom numbering.

Two descriptors are both named \texttt{BalabanJ} and are not the same quantity --- RDKit's uses a
bond-order-weighted distance matrix, Mordred's passes RDKit's own formula an unweighted one, and
on naphthalene they are 2.888052 and 1.925368 --- so both are emitted under distinct names.
RDKit's $\chi$ is internally inconsistent in the same way: \texttt{Chi0n}/\texttt{Chi1n} count
all atoms and bonds while \texttt{Chi0v}/\texttt{Chi1v} are heavy-atom only, so
\texttt{Chi0n}\,$-$\,\texttt{Chi0v} is exactly the explicit-hydrogen count while the variants
agree identically at $k \geq 2$, and its path finder counts rings as paths by its own documented
design. We match RDKit throughout, which makes explicit hydrogens value-changing rather than
cosmetic ($d_3$-ethanol has $\chi_{2n} = 0.158114$ against ethanol's $0.316228$) and is why no
hydrogen stripping is applied. Crippen and E-state typing dominated cost, being a generic
subgraph search answering a question that for most rows concerns an atom and its immediate
neighbours; both are replaced by a dispatch on (element, aromaticity) and a per-row predicate
over the atom's own properties, with branch specifications resolved by bounded backtracking as a
\emph{system of distinct representatives} rather than a tally of bond types, since SMARTS maps
distinct query atoms to distinct molecule atoms and the two readings diverge as soon as two
branch specifications overlap. Multi-typing is preserved, since an atom legitimately receives a
tuple of type names rather than a best match. This takes Crippen from roughly 78 to
1.53~$\mu$s/molecule and E-state from 636 to 0.834~$\mu$s, bit-exactly.

Two numerical kernels are written out rather than called from BLAS/LAPACK, because identical
source measured 138.09 against 218.93~$\mu$s/molecule depending purely on which BLAS the host
supplied, and a released package cannot have its values decided by the machine it was built on.
BCUT2D's extremal eigenvalues use Householder tridiagonalization followed by a
Pal--Walker--Kahan QL/QR sweep, precisely the two stages \texttt{dsytd2} and \texttt{dsterf}
perform and in the same order, so agreement is structural rather than tuned --- maximum deviation
$4.26 \times 10^{-13}$ over all 395{,}620 Burden matrices --- and faster than the vendor library
at 27 heavy atoms, where a tuned library is mostly wrapper. The resistance solve is a reference
LU factorization with partial pivoting rather than the shorter, better-conditioned Cholesky that
$L + J/k$ admits, because the reference is \texttt{numpy.linalg.inv}, which is LU with partial
pivoting, and matching the algorithm is what makes the arithmetic identical rather than merely
close; bit-identity is the binding constraint because resistance bin edges sit at $0.1$, $0.5$,
$1.0$ and $2.0$ and real molecular values land exactly on them --- on tetra-\emph{tert}-butyl
tetrahedrane every atom pair has $\Omega = 0.5$ --- so values within $10^{-9}$ of an edge are
snapped before binning and all three solvers then produce identical bin vectors on all 98{,}905
molecules. Sentinels are propagated rather than tidied away: RDKit writes $10^8$ for unreachable
pairs and \texttt{BalabanJ} sums those rows, and skipping them instead gives a finite, plausible
and wrong answer on 16\% of an adversarial corpus but on roughly one molecule of a drug-like one.
Verification is per-decision rather than per-column, since a count such as \texttt{NsCH3} lets
two opposite mistypings cancel, so we check the atom's type tuple and the Crippen row index over
2{,}868{,}290 and 2{,}869{,}048 atoms respectively, and 98{,}905 molecules for the block columns,
against RDKit 2025.9.2 and Mordred 1.2.0. Exactness is version-relative: the E-state and Crippen
data are byte-identical between RDKit 2025.9.2 and 2026.3.5 and our C++ scores 100\% against
both, yet 19 of 100{,}000 molecules receive different E-state atom types because furans fused
into large macrocycles stop being perceived as aromatic, so the RDKit version is part of the
released column specification.

% ---------------------------------------------------------------------------------------------
\subsection{Novel descriptors}
% ---------------------------------------------------------------------------------------------

We add four blocks, 140 columns, that no descriptor in either library computes. The generating
question is deliberately not ``what information is ECFP missing?'' --- by the data-processing
inequality, and because ECFP is near-injective on real molecules, the answer is essentially
nothing. It is instead: \emph{what is present in the fingerprint but not in a form a
gradient-boosted tree can construct?} Three shapes fail that test. Pairwise sums
$\sum_{ij} f(i,j)$ need quadratic interaction across 2048 hashed bits; ratios need division,
which trees cannot do; and cyclic patterns are outside colour refinement altogether, since
WL-1 distinguishes exactly the graphs that homomorphism counts \emph{from trees} distinguish.
Only the third is an expressivity limit in the theorem sense --- the other two are conditioning
--- and the blocks are shaped accordingly.

\paragraph{Resistance (60 columns).} Treat every bond as a unit resistor and let $\Omega_{ij}$
be the effective resistance between atoms $i$ and $j$, computed from the Moore--Penrose
pseudo-inverse of the graph Laplacian, taken per connected component. The block is built around
$\Delta_{ij} = d_{ij} - \Omega_{ij}$, the gap between shortest-path and resistance distance.
On a tree there is exactly one route between any two atoms, so $\Omega = d$ and $\Delta \equiv 0$:
the quantity measures path multiplicity --- ring fusion --- and \emph{cannot} proxy for size,
weight or atom count. We emit centred autocorrelations of four atomic properties binned by
$\Delta$, the Kirchhoff index and its normalisations, and pooled random-walk return
probabilities $\mathrm{diag}(S^k)$ for $k \in \{2,3,4,6,8,12,16\}$ with
$S = D^{-1/2} A D^{-1/2}$. To be precise about the orthogonality claim: the 30 $\Delta$-derived
columns are identically zero on acyclic molecules, while the Kirchhoff and random-walk columns
are not (on a tree $K_f$ equals the Wiener index).

\paragraph{Cycles (33 columns).} Exact counts of simple cycles of length 3 to 8, per-atom
participation, and hetero/aromatic ring typing. This is the block that targets the WL-1 gap
directly. It is also not \texttt{RingCount}: the SSSR is a \emph{basis} for the cycle space, not
a list of cycles, and the two diverge sharply --- cubane's cycle space has dimension 5, yet the
molecule contains 16 distinct six-membered cycles and 28 cycles of length $\leq 8$ in total.
Lengths 3--5 are obtained in closed form from traces of $A^3$, $A^4$, $A^5$; 6--8 by bounded
depth-first search from each cycle's lowest-indexed vertex.

\paragraph{Conjugation (24 columns).} Union-find over RDKit's conjugated-bond flags recovers the
$\pi$-system connected components, then their count, sizes, size histogram, topological
diameter, branch points (cross-conjugation), heteroatom content and non-aromatic extent.
Coverage of conjugated-system \emph{topology} in both libraries is zero. The motivating case is
that anthracene and a C$_{14}$ polyene are both a single conjugated system of 14 atoms and are
separated only by shape: linearity (diameter over size) is 0.54 against 1.00.

\paragraph{Stereo (23 columns).} Signed CIP parities $s_i \in \{-1,0,+1\}$, combined at two
orders. Odd-order terms such as $\sum_i s_i$ flip sign under reflection and therefore separate
enantiomers; even-order terms $\sum_{d(i,j)=k} s_i s_j$ are mirror-invariant and separate
diastereomers. Analogous terms and cross terms are built on E/Z bond parities. The motivation
is a measurement: the descriptor union is \emph{completely} stereo-blind --- zero of RDKit's 217
and zero of Mordred's 1{,}613 descriptors change value across any enantiomer pair. This block is
therefore conditioning rather than new information, since ECFP with chirality enabled already
carries stereo in 4--10 bits, and we report it as such.

Cycles, conjugation and stereo each return exactly zero when their axis is absent: an acyclic
molecule gets 33 zeros, a molecule with no conjugated bond gets 24, a molecule with no
stereocentre and no E/Z bond gets 23. Resistance is zero only on its 30 $\Delta$-derived
columns, as noted above. That is orthogonality by construction rather than by hope.

\paragraph{Two implementation details that changed values.} First, \emph{a cycle is not its
vertex set}. In $K_4$ three distinct 4-cycles share a single vertex set, so de-duplicating
enumerated cycles on their sorted vertex tuple silently merges two of the three. On a
tetrahedrane derivative this removed exactly $0.125$ from both $\chi_{4n}$ and $\chi_{4v}$
(8.787913 against RDKit's 8.912913). Cycles are canonicalised on
$\mathrm{path}[1] < \mathrm{path}[-1]$ instead. Second, $\Omega$ is a rational function of the
graph, so $\Delta$ lands \emph{exactly} on bin edges for symmetric molecules, and the binned
columns are then decided by floating-point noise: two LAPACK implementations on the same machine
disagreed on 9.00\% of a 98{,}905-molecule corpus. Values within $10^{-9}$ of an edge are snapped
to it before binning, which changed 13.17\% of molecules and makes the columns
implementation-independent.

% ---------------------------------------------------------------------------------------------
\subsection{XGBoost implementation}
% ---------------------------------------------------------------------------------------------

% AUTHOR NOTE -- resolve before submission. The repo currently contains two XGBoost
% configurations: vendor/chemtfm/models/xgb.py (used by the LOCKED-suite scripts) sets
% subsample=0.8 and colsample_bytree=0.8; gpu/dev_grid_v2.py does NOT set either. The two
% also use different scaffold_folds implementations. The text below describes the vendored
% config. Either unify them or say which result used which -- "the same XGBoost throughout"
% is not currently true.

The downstream head is a deliberately \emph{untuned} XGBoost: 300 rounds, \texttt{max\_depth} 6,
$\eta = 0.1$, \texttt{subsample} 0.8, \texttt{colsample\_bytree} 0.8,
\texttt{min\_child\_weight} 1, \texttt{tree\_method="hist"}. No hyperparameter search, no early
stopping, no validation set. The head is a measuring instrument for the representation, and
per-dataset tuning would move the bar between arms and make the comparison unreproducible.
Missing values are passed to XGBoost unimputed --- NaN is a first-class input, and the learned
default split direction is more informative than a median. Fingerprint counts enter as
$\log(1+x)$, since raw counts are heavy-tailed. No feature standardisation is applied to the
model input.

Evaluation is 5-fold Bemis--Murcko \textbf{scaffold} cross-validation, never random: random
splits put near-identical analogues on both sides of the split. Folds are computed once per
dataset and shared by every arm, so fold difficulty cancels in the paired differences. Results
are aggregated by mean rank and win-rate rather than by averaging raw RMSE across datasets
whose scales differ by two- to three-fold.

\paragraph{One extra hyperparameter is required, and it is not a standard one.} XGBoost's
\texttt{colsample\_by*} samples columns \emph{uniformly}. That is the wrong prior for a matrix
that is half sparse fingerprint bits and half dense descriptors: a descriptor takes an
informative value on nearly every molecule, while a given ECFP bit is set in roughly 2\% of
them, so descriptors win the column-sampling lottery on \emph{availability} rather than on merit
for the task. We therefore set XGBoost's \texttt{feature\_weights} to a single scalar $w$ on the
2048 fingerprint columns and $1$ on the descriptors, which converts ``descriptors on or off''
into a continuous knob. Two implementation details matter: \texttt{feature\_weights} does
nothing unless \texttt{colsample\_by*} $< 1$, and it must be applied through
\texttt{colsample\_bynode} (we use $0.3$) rather than \texttt{colsample\_bytree}, because
weights act at each sampling event and per-node sampling gives the reweighting far more chances
to bite.

Measured over 34 datasets under scaffold folds (RMSE, lower is better):

\begin{center}
\begin{tabular}{lrrrr}
\toprule
arm & ACE all & ACE cliff & ACE non-cliff & MoleculeNet \\
\midrule
fingerprint only     & 0.7854 & 0.8610 & 0.7338 & 0.9487 \\
descriptors only     & 0.8460 & 0.9065 & 0.8055 & 0.9358 \\
$w = 1$ (XGBoost default) & 0.8080 & 0.8752 & 0.7626 & 0.8901 \\
$w = 10$             & 0.7882 & 0.8605 & 0.7389 & \textbf{0.8712} \\
$w = 100$            & \textbf{0.7721} & 0.8464 & 0.7214 & 0.8931 \\
\bottomrule
\end{tabular}
\end{center}

The XGBoost default, $w = 1$, is \emph{dominated on both suites}. We recommend tuning $w$ in the
inner cross-validation loop alongside the other head hyperparameters; if it is not tuned,
$w = 10$ rather than $1$. The two curves have different shapes, which is the substantive result:
the potency suite improves monotonically out to $w = 100$, where the reweighted model beats a
pure fingerprint, whereas the physicochemical suite has an interior optimum near $w \approx 10$
and degrades by $w = 100$.

\paragraph{The mechanism is not the one we tested for.} The sweep was run to test the
hypothesis that descriptors cannot resolve \emph{activity cliffs}, being smooth functions of
structure, so up-weighting fingerprints should help specifically on cliff compounds. That
prediction failed: going from $w = 1$ to $w = 100$ improves non-cliff RMSE by 5.4\% and cliff
RMSE by only 3.3\%, so the gain is \emph{larger} on the molecules that are not cliffs. The
surviving explanation is a task-family effect rather than a cliff effect --- MoleculeACE
datasets are congeneric series, so descriptors vary little across the whole dataset, not merely
across cliff pairs.

% AUTHOR NOTE: the sweep above is n_repeats=1, fold_seed=0, xgb_seed=0 -- a single seed with no
% error bars, and several gaps are 0.01-0.04. The 3.3%-vs-5.4% cliff comparison in particular is
% exactly the size that needs replicates. Re-run multi-seed before this table goes in.
